\section*{第 \ref{ch:b1:fluid-statics} 章——流体静力学}

\begin{solution}{exo:b1:fluid-statics:1}
$P = P_0 + \rho gh$: \qty{2.0e5}{Pa} (\qty{2.0}{bar}), \qty{1.1e6}{Pa}
（\qty{11}{bar}），\qty{1.0e7}{Pa}（\qty{101}{bar}）。在 \qty{300}{m} 深处：$P =
\qty{3.1e6}{Pa}$，作用在 \qty{1}{m^2} 上的力为 \qty{3.1}{MN}（净值，
扣除内部 \qty{1}{atm} 的反压后为 \qty{3.0}{MN}）。
\end{solution}

\begin{solution}{exo:b1:fluid-statics:2}
$h = P/\rho g$：汞柱 \qty{0.760}{m}；水柱 \qty{10.3}{m}。水泵
至多能在水面上方造出真空，此时大气
压将水柱托起 \qty{10.3}{m}，不能再高。
\end{solution}

\begin{solution}{exo:b1:fluid-statics:3}
$F_2 = F_1S_2/S_1 = \qty{10}{kN}$（相当于一吨）；$d_2 = d_1S_1/S_2 = \qty{2.0}{mm}$
（排开的体积相同，做功相同）。
\end{solution}

\begin{solution}{exo:b1:fluid-statics:4}
两臂中油–水界面所在水平面处的压强相同：$\rho_og
h_o = \rho_wgh_w$，$h_w = 0.85 \times 12 = \qty{10.2}{cm}$，即另一臂中该水平面以上
的水柱高度；油面
则高出 $12 - 10.2 = \qty{1.8}{cm}$。
\end{solution}

\begin{solution}{exo:b1:fluid-statics:5}
$H = RT/Mg = 8.314 \times 273/(0.029 \times 9.81) = \qty{8.0}{km}$; $P/P_0 =
\eu^{-z/H}$：在 \qty{5}{km} 处为 $0.54$，在 \qty{8848}{m} 处为 $0.33$；压强减半的高度为 $H\ln2 =
\qty{5.5}{km}$——即在\cref{ch:b1:kinetic-theory}中所用的数值。
\end{solution}

\begin{solution}{exo:b1:fluid-statics:6}
$F = \tfrac12\rho gh^2L = \qty{1.96e8}{N}$；\omterm{prop:b1:fluid-statics:wall}{压力中心}位于 \qty{13.3}{m}
深处，即底部以上 \qty{6.7}{m}；对底部的力矩为 $F \times
6.7 = \qty{1.3e9}{N.m}$。将坝体朝水的一侧弯成拱形，使推力
转化为拱的压缩，并传至山谷两侧。
\end{solution}

\begin{solution}{exo:b1:fluid-statics:7}
$917/1025 = 89\%$；原木 $60\%$。筏子：完全浸没时的浮力 $1000 \times
2.0 \times 9.81 = \qty{19.6}{kN}$，自重 \qty{11.8}{kN}，
尚余 \qty{7.8}{kN}，可载十人\dots{}第十人脚湿：九人可保持干燥
(\qty{800}{kg}).
\end{solution}

\begin{solution}{exo:b1:fluid-statics:8}
浸没体积为 $m/\rho$：水中 \qty{20}{cm^3}；盐水中 $20/1.1 = \qty{18.2}{cm^3}$，
少 $\qty{1.8}{cm^3}$，即杆身多出 $1.8/0.50 = \qty{3.6}{cm}$，
故露出 \qty{7.6}{cm}。灵敏度 $\approx m/(\rho^2S) \approx \qty{40}{cm}$ 每单位
相对密度（相对密度变化 $0.1$ 时约 \qty{3.6}{cm}）。
\end{solution}

\begin{solution}{exo:b1:fluid-statics:9}
排开体积 $2800/1.025 = \qty{2732}{m^3}$：\qty{268}{m^3}（$9\%$）露出
水面。中性浮力时质量为 $1025 \times 3000 = \qty{3075}{t}$：需注入 \qty{275}{t}。
淡水中：浮力降至 \qty{3000}{t}，超重 \qty{75}{t}——将下沉，
除非抽出 \qty{75}{t}。
\end{solution}

\begin{solution}{exo:b1:fluid-statics:10}
排开空气：$10^{-2} \times 1.20 = \qty{12.0}{g}$；内部氦气 $12.0 \times
4/29 = \qty{1.7}{g}$；橡胶 \qty{3.0}{g}，净升力 \qty{7.3}{g}——可承托七米
长的轻绳。
\end{solution}

\begin{solution}{exo:b1:fluid-statics:11}
力矩 $\int_0^h\rho gx\cdot xL\,\dd x = \rho gLh^3/3$，压力 $\rho gLh^2/2$：
\omterm{def:b1:angular-momentum:moment}{力臂} $2h/3$。自 $h_1$ 至 $h_2$：力矩 $\rho gL(h_2^3 - h_1^3)/3$，压力
$\rho gL(h_2^2 - h_1^2)/2$：深度 $\frac23(h_2^3 - h_1^3)/(h_2^2 - h_1^2)$。
\end{solution}

\begin{solution}{exo:b1:fluid-statics:12}
多浸没 $x$ 时，附加浮力 $\rho gSx$ 竖直向上：$m\ddot x = -\rho gSx$，
$\omega_0^2 = \rho gS/m$, $T = 2\pi\sqrt{m/\rho gS} = 2\pi\sqrt{0.020/(1000 \times
9.81 \times 5\times10^{-5})} = \qty{1.3}{s}$.
\end{solution}

\begin{solution}{pb:b1:fluid-statics:1}
\textbf{1.} $P = 1.013\times10^5 + 1025 \times 9.81 \times 300 = \qty{3.12e6}{Pa}
= \qty{31}{bar}$.

\textbf{2.} $S = \pi(0.40)^2 = \qty{0.503}{m^2}$：合力 $(P - P_0)S =
\qty{1.52e6}{N}$——相当于 155 吨，即一辆满载卡车压在舱盖上。

\textbf{3.} 龙骨再深 \qty{10}{m}：压强增加 $+\rho g \times 10 = \qty{1.0e5}{Pa}$，
约 $3\%$：船体受载近似均匀。

\textbf{4.} $\Delta\rho/\rho = \chi_T\Delta P = 4.6\times10^{-10} \times 3.0\times10^6
= 0.14\%$：对压强可忽略（仅几千帕，占
三百万帕中的一部分）。

\textbf{5.} 外部三十巴，内部一巴：船体是一种承受向内压力的压力容器，
圆柱（或球面）壳能把压力转化为
沿壁的压缩，钢材对此抵抗力强，而平板
则会弯曲。

\textbf{6.} $\Delta P\,R/e = 3.0\times10^6 \times 5.0/0.030 = \qty{5.0e8}{Pa}
= \qty{500}{MPa}$：接近屈服应力——\qty{300}{m} 已接近
此船体的极限，更深的船需要更厚或更强的钢材
（或钛）。

\textbf{7.} 排开体积 $m_0/\rho = 2800/1.025 = \qty{2732}{m^3}$：\qty{268}{m^3}
露出水面（$9\%$）。

\textbf{8.} 中性浮力：$m = \rho V = \qty{3075}{t}$：注入 \qty{275}{t}。

\textbf{9.} 浮力减少 $0.1\%$ 的 $\rho Vg$：$0.001 \times 3075 \times 9.81
\times10^3 = \qty{30}{kN}$，方向向下（相当于三吨）。更深 $\Rightarrow$ 压缩增加
$\Rightarrow$ 变重：不稳定——潜艇必须持续调整配平
（钢质艇体的压缩比水小，这有所帮助；
艇体过于柔软则会沉没而无法回升）。

\textbf{10.} 浮力 $1000 \times 3000 \times 9.81 = \qty{29.4}{MN}$，与
重力 $3075 \times 9.81 \times 10^3 = \qty{30.2}{MN}$：向下 \qty{0.74}{MN}
（\qty{75}{t}）；抽出 \qty{75}{t}。

\textbf{11.} 立即偏轻 \qty{20}{t}：需注入 \qty{20}{t} 水
（\qty{19.5}{m^3}）至平衡水舱。

\textbf{12.} 水只能被压强高于外部海水的物体排出；
瓶中压缩空气将水压出——仅当
瓶内压强超过该深度处环境压强
时才能做到。

\textbf{13.} $PV$ 为常量：\qty{275}{m^3} 在 \qty{31}{bar} 下需 $275 \times 31
= \qty{8500}{m^3}$ 的 \qty{1}{atm} 空气——在 \qty{200}{bar} 下贮存于
\qty{43}{m^3} 的气瓶中。

\textbf{14.} \qty{10}{m} 深处为 \qty{2.0}{bar}；\qty{30}{m} 深处为 \qty{4.0}{bar}。

\textbf{15.} $PV$ 为常量：$6.0/4.0 = \qty{1.5}{L}$——胸廓
向内压缩；自由潜水者对此深有体会。

\textbf{16.} 在 \qty{4}{bar} 下吸入的空气膨胀至四倍，即 \qty{24}{L}，
置于仅能容纳 \qty{6}{L} 的肺中：会导致破裂。上升时切勿屏住呼吸，
应持续呼气。

\textbf{17.} $\Delta P = \rho g \times 1.0 = \qty{1.0e4}{Pa}$；受力为 $\qty{1.0e4}{Pa}
\times 0.05 = \qty{500}{N}$，作用于胸部：呼吸肌无法
托起五十千克的重物——呼吸管只在最初几分米内有效。

\textbf{18.} $12 \times 200 = \qty{2400}{L}$（\qty{1}{bar} 下），即
在 \qty{4}{bar} 下为 \qty{600}{L}：\qty{30}{min}。

\textbf{19.} $V = 10^8/1025 = \qty{97600}{m^3}$。

\textbf{20.} 在淡水中须排开 $10^5\ \unit{m^3}$：还需多排开 \qty{2400}{m^3}，
即吃水加深 $2400/8000 = \qty{0.30}{m}$。

\textbf{21.} 还需多排开 $10^4/1.025 = \qty{9756}{m^3}$：吃水加深 \qty{1.2}{m}。

\textbf{22.} 空载时，船的 $G$ 位置较高，浸没的船体很少：
\omterm{rem:b1:fluid-statics:stability}{稳心}可能低于 $G$，船在涌浪中可能倾覆。海
水若低置于船体内，会降低 $G$ 并加大吃水。

\textbf{23.} $\rho gh = \qty{2.0e5}{Pa}$，即 \qty{0.2}{MN/m^2}：约为
潜艇所受 \qty{3}{MN/m^2} 的十五分之一。

\textbf{24.} 否：潜艇所在深度的压强为 $P_0 + \rho gh$，
仅由深度决定。油轮的重力由其所排开的水承担，
这使海平面处处升高一个极小的量，
没有额外载荷传到潜艇上。

\textbf{25.} $\dd P/\dd z = -\rho g$——由此得 $P = P_0 + \rho gh$，
对此处每一深度、受力与肺容积均成立——以及阿基米德定理，它是
该关系在闭合曲面上的积分。
\end{solution}
